\chapter{Los números naturales y enteros}
\label{cap:naturales-enteros}

\section{Aritmética y orden en los naturales}
\label{sec:aritmetica-naturales}
Se adopta $\N=\{0,1,2,\ldots\}$ y $\N^*=\N\setminus\{0\}$.
Los axiomas de Peano y la inducción se recogen en la
sección~\ref{sec:induccion}. El sucesor es $s(n)=n+1$; cada natural no
nulo tiene un único predecesor y ningún natural coincide con su sucesor.

La suma y el producto se definen por recurrencia en el segundo argumento:
\begin{align*}
 m+0&=m,&m+s(n)&=s(m+n),\\
 m\cdot0&=0,&m\cdot s(n)&=mn+m.
\end{align*}
La inducción permite trasladar estas definiciones a todos los naturales
y demostrar las reglas aritméticas:
\begin{itemize}
\item Suma y producto son asociativos y conmutativos; tienen neutros
$0$ y $1$. El producto distribuye respecto a la suma y $m0=0$.
\item $m+p=n+p\imp m=n$. Si $p\ne0$, $mp=np\imp m=n$.
\item $m+n=0\eqv m=n=0$; $mn=0\eqv m=0$ o $n=0$.
\end{itemize}
Para $a\ne0$, $a^0=1$ y $a^{n+1}=a^na$; para $n>0$, $0^n=0$.
No se fija aquí un valor aritmético general para $0^0$.
Si $a,b\in\N^*$ y $m,n\in\N$,
\[
 a^ma^n=a^{m+n},\qquad (a^m)^n=a^{mn},\qquad (ab)^n=a^nb^n.
\]

El orden se obtiene de la suma:
\[
 m\le n\eqv \exists p\in\N: m+p=n,
 \qquad m<n\eqv m+1\le n.
\]
Es total, tiene mínimo $0$ y no tiene máximo. Sumar un natural o
multiplicar por un natural conserva las desigualdades; multiplicar por
un natural positivo conserva también las estrictas.
Es un orden \textbf{discreto}: no hay naturales entre $n$ y $n+1$.
Además, $\N$ está \textbf{bien ordenado}: todo subconjunto no vacío tiene
mínimo. Un subconjunto no vacío acotado superiormente tiene máximo.

\section{Conjuntos finitos y recuentos}
\label{sec:recuentos}
Un conjunto es \textbf{finito} si es vacío o es equipotente a
$\{1,\ldots,n\}$ para algún $n\ge1$. Ese $n$ es único y es su cardinal;
el vacío tiene cardinal $0$. Una inyección entre conjuntos finitos
obliga a que el cardinal del origen no supere al del destino.
En particular, $\N$ es infinito.

Todo subconjunto de un conjunto finito es finito; si es propio, su
cardinal es estrictamente menor. Para subconjuntos de $\N$, ser finito
equivale a estar acotado superiormente. La imagen de un conjunto finito
$A$ por $f$ tiene cardinal a lo sumo $|A|$, con igualdad si y solo si
$f$ es inyectiva. Entre conjuntos finitos del mismo cardinal,
inyectividad, sobreyectividad y biyectividad son equivalentes
(sección~\ref{sec:cardinalidad}).

Para conjuntos finitos $A,B$,
\[
 |A\cup B|+|A\cap B|=|A|+|B|,\qquad |A\times B|=|A||B|.
\]
Si son disjuntos, el cardinal de la unión es la suma. Los recuentos
fundamentales, con $|A|=n$ y $|B|=m$, son:
\begin{center}
\begin{tabular}{@{}ll@{}}
\toprule
Objeto & Número\\
\midrule
Aplicaciones $A\to B$, si $m>0$ & $m^n$\\
Inyecciones $A\to B$, si $0\le n\le m$ & $\dfrac{m!}{(m-n)!}$\\[6pt]
Biyecciones $A\to B$, si $m=n$ & $n!$\\
Subconjuntos de $A$ & $2^n$\\
Subconjuntos de $n$ elementos de $B$ & $\displaystyle\binom mn=\dfrac{m!}{n!(m-n)!}$\\
\bottomrule
\end{tabular}
\end{center}
Las inyecciones cuentan \textbf{variaciones} sin repetición; las
biyecciones de un conjunto consigo mismo, \textbf{permutaciones};
los subconjuntos de tamaño fijado, \textbf{combinaciones}.
Se toma $0!=1$. Si $n>m$ no hay inyecciones ni subconjuntos de $n$
elementos de $B$. Si $B=\varnothing$, hay una aplicación desde el vacío
y ninguna desde un conjunto no vacío. Para $0\le n\le m$,
$\binom mn=\binom m{m-n}$; la identidad de Pascal es
\[
 \binom{m+1}{n}=\binom m{n-1}+\binom mn\qquad(1\le n\le m).
\]
Se explica separando los subconjuntos según contengan o no un elemento
fijado. El número $\binom mn$ no cuenta ordenaciones: el orden de los
elementos del subconjunto es irrelevante.

\section{Infinitud y numerabilidad}
\label{sec:numerabilidad}
Un conjunto infinito puede ser equipotente a una parte propia:
$n\mapsto n+1$ identifica $\N$ con $\N^*$.
Un conjunto \textbf{numerable} es equipotente a $\N$; en esta terminología
los conjuntos finitos no son numerables.

\textbf{Teorema de Cantor.} Para cualquier conjunto $A$,
\[
 \operatorname{card}A<\operatorname{card}\mathcal P(A).
\]
La aplicación $a\mapsto\{a\}$ es inyectiva. Ninguna $h:A\to\mathcal P(A)$
es sobreyectiva: el conjunto $D=\{a\in A:a\notin h(a)\}$ no puede ser
$h(d)$ para ningún $d$, pues $d\in D\eqv d\notin D$.
Así se obtienen cardinales infinitos cada vez mayores tomando conjuntos
de partes sucesivamente.

Todo subconjunto infinito de $\N$ es numerable: se enumera tomando su
mínimo y después el mínimo de lo que queda. No queda ningún elemento
fuera de la lista, porque por debajo de cada natural solo hay un número
finito de naturales. Transportando este resultado por una biyección,
todo subconjunto de un conjunto numerable es finito o numerable.

La enumeración de $\N^2$ por diagonales $n+m=0,1,2,\ldots$ muestra que
$\N^2$ es numerable. Se deducen las siguientes reglas:
\begin{itemize}
\item El producto de uno o más conjuntos numerables, en número \emph{finito},
es numerable.
\item La unión de una familia numerable de conjuntos numerables es numerable:
se enumeran sus elementos mediante pares de índices y se suprimen repeticiones.
La elección simultánea de enumeraciones se hace bajo el axioma de elección.
\item Una unión numerable de conjuntos finitos o numerables es finita o
numerable; si contiene un subconjunto infinito, es numerable.
\end{itemize}
No se extiende la primera regla a productos de infinitos factores.
Los conjuntos $\Z$ y $\Q$ son numerables: los enteros se listan alternando
signos, y los racionales se representan por pares de enteros con
denominador positivo. También es numerable el conjunto de subconjuntos
finitos de $\N$, separándolos según su número de elementos.

\section{Construcción y orden de los enteros}
\label{sec:enteros}
La construcción por clases de equivalencia de la
sección~\ref{sec:equivalencia} utiliza
\[
 (a,b)\sim(c,d)\eqv a+d=b+c\qquad(a,b,c,d\in\N).
\]
La clase $[a,b]$ representa la diferencia $a-b$; las operaciones son
\[
 [a,b]+[c,d]=[a+c,b+d],\qquad
 [a,b][c,d]=[ac+bd,ad+bc].
\]
No dependen de los representantes. El cero es $[0,0]$, la unidad es
$[1,0]$ y el opuesto de $[a,b]$ es $[b,a]$. Se obtiene un anillo
conmutativo unitario sin divisores de cero.
La identificación $n\mapsto[n,0]$ incluye $\N$ en $\Z$ y conserva sus
operaciones.

Cada clase tiene un representante de la forma $[m,0]$ o $[0,n]$.
Los no negativos son las clases $[m,0]$ y el orden se define por
\[
 \alpha\le\beta\eqv\beta-\alpha\ge0.
\]
Así $\Z$ es un anillo totalmente ordenado y la inclusión de $\N$
conserva el orden. Todo subconjunto no vacío de $\Z$ acotado inferiormente
tiene mínimo; si está acotado superiormente, tiene máximo.
$\Z$ no está bien ordenado y no tiene mínimo ni máximo.
Es \textbf{arquimediano}: dados $\alpha>0$ y $\beta\in\Z$, existe
$n\in\N$ con $n\alpha>\beta$.

\section{División entera y subgrupos de los enteros}
\label{sec:division-entera}
Para $a,b\in\Z$ con $b\ne0$, existen únicos $q,r\in\Z$ tales que
\[
 a=qb+r,\qquad 0\le r<|b|.
\]
$q$ es el cociente y $r$ el resto de la \textbf{división entera}.
La cota del resto da la unicidad: la diferencia de dos restos no puede
ser un múltiplo no nulo de $b$.
Con dividendo negativo no basta truncar el cociente hacia cero:
$-17=(-4)\cdot5+3$ tiene resto admisible, mientras que $-2$ no lo sería.

Todo subgrupo de $(\Z,+)$ es $n\Z=\{nk:k\in\Z\}$ para un único
$n\in\N$. Si el subgrupo no es $\{0\}$, se toma su menor elemento
positivo $n$. Al dividir cualquier elemento por $n$, el resto pertenece
al subgrupo; su minimalidad obliga a que sea cero.
Todos estos subgrupos son ideales de $\Z$, por lo que \textbf{todo ideal
de $\Z$ es principal}. El caso $n=0$ corresponde a $\{0\}$.

\section{Divisibilidad, máximo común divisor y mínimo común múltiplo}
\label{sec:divisibilidad}
Para $a,b\in\Z$,
\[
 b\mid a\eqv\exists q\in\Z: a=qb
 \eqv a\in b\Z\eqv a\Z\subseteq b\Z.
\]
Todo entero divide a $0$; $0$ solo divide a $0$. Los enteros $1$ y $-1$
dividen a todos los enteros. La divisibilidad es un orden en $\N^*$,
pero no en $\Z$: $a$ y $-a$ se dividen mutuamente.

Para $a,b>0$, se definen los enteros positivos $d,m$ por
\[
 a\Z+b\Z=d\Z,\qquad a\Z\cap b\Z=m\Z,
 \quad a\Z+b\Z=\{au+bv:u,v\in\Z\}.
\]
Son, respectivamente, el \textbf{máximo común divisor} y el
\textbf{mínimo común múltiplo}:
\begin{itemize}
\item $d=\gcd(a,b)$ divide a $a,b$, y todo divisor común divide a $d$.
Es además la menor combinación lineal positiva $au+bv$.
\item $m=\operatorname{mcm}(a,b)$ es múltiplo de $a,b$ y divide a todo
múltiplo común. Es el menor múltiplo común positivo.
\end{itemize}
La \textbf{identidad de Bézout} afirma que existen $u,v\in\Z$ tales que
\[
 d=au+bv.
\]
Para enteros no nulos se usa $\gcd(a,b)=\gcd(|a|,|b|)$.
Se llaman \textbf{primos entre sí} si este número es $1$;
equivalentemente, $au+bv=1$ para ciertos enteros $u,v$.

\textbf{Algoritmo de Euclides.} Para calcular el máximo común divisor
de $a\ge b>0$, dividir y reemplazar el par $(a,b)$ por $(b,r)$ mientras
$r\ne0$. La igualdad $a=qb+r$ implica
$a\Z+b\Z=b\Z+r\Z$, así que el máximo común divisor se conserva.
Los restos positivos decrecen y el último resto no nulo es $d$;
si la primera división es exacta, $d=b$. Sustituir hacia atrás las
igualdades de las divisiones expresa $d$ como $au+bv$.

\textbf{Teorema de Gauss.} Si $a,b\ne0$, $\gcd(a,b)=1$ y $a\mid bc$,
entonces $a\mid c$. En efecto, multiplicar $au+bv=1$ por $c$ expresa
$c$ como suma de dos múltiplos de $a$.

\section{Naturales y cardinales finitos}
Los naturales pueden interpretarse como los cardinales de los conjuntos
finitos. La suma cuenta uniones disjuntas y el producto cuenta pares
ordenados. Añadir un elemento nuevo realiza el sucesor $n\mapsto n+1$;
esta interpretación enlaza la aritmética con los recuentos de la
sección~\ref{sec:recuentos}.
